\documentclass[conference]{IEEEtran}
\IEEEoverridecommandlockouts
\usepackage{cite}
\usepackage{amsmath,amssymb,amsfonts}
\usepackage{algorithmic}
\usepackage{graphicx}
\usepackage{textcomp}
\usepackage{xcolor}
\usepackage{booktabs}
\usepackage{multirow}
\usepackage{url}
\usepackage{hyperref}

\def\BibTeX{{\rm B\kern-.05em{\sc i\kern-.025em b}\kern-.08em
    T\kern-.1667em\lower.7ex\hbox{E}\kern-.125emX}}

\begin{document}

\title{JanAwaaz AI: A Privacy-Preserving Autonomous Multi-Agent Framework for Zero-Knowledge Public Grievance Redressal and Real-Time Civic Triage\\
\thanks{This research represents original engineering and architectural work on the JanAwaaz AI anonymous public grievance ecosystem, deployed as an open-access civic technology initiative.}
}

\author{\IEEEauthorblockN{Prakhar}
\IEEEauthorblockA{\textit{Department of Computer Science and Engineering} \\
\textit{Independent Research Initiative}\\
Email: wizardmathsprakhar08@gmail.com}
}

\maketitle

\begin{abstract}
Public grievance redressal mechanisms are the foundational pillar of participatory democracy and municipal governance. However, contemporary civic portals suffer from severe citizen attrition stemming from two systemic pathologies: the pervasive fear of administrative retaliation due to non-anonymized personal disclosures, and catastrophic operational bottlenecks caused by manual triage, duplicate floodings, and uncalibrated urgency assessments. In this paper, we propose \textbf{JanAwaaz AI}, an end-to-end, privacy-preserving, autonomous multi-agent framework designed to facilitate zero-knowledge citizen grievance reporting and real-time administrative triage. JanAwaaz AI decomposes civic issue processing into a decoupled five-agent sequential pipeline: (1) an Autonomous Ingestion and Zero-Knowledge PII Redaction Agent that scrubs phone numbers, email identifiers, Aadhaar numbers, and named entities in linear time while issuing cryptographically collision-resistant tracking tokens; (2) a Semantic Categorization Agent executing lexical-domain intent disambiguation across civic departments; (3) a Multi-Factor Heuristic Urgency Scoring Engine evaluating cognitive panic, lexical severity, and structural threats; (4) a Vector-Space Duplicate Detection Agent implementing TF-IDF cosine metric clustering to eliminate redundant ticket processing; and (5) an Automated Routing and Cryptographic Lifecycle State Machine governing officer dispute queues. Deployed over a distributed cloud microservice architecture with progressive web (PWA) and React Native mobile clients, JanAwaaz AI was evaluated against a multi-department benchmark corpus. Empirical results demonstrate a 99.4\% PII redaction F1-score, a 91.8\% departmental routing accuracy, a 64\% reduction in duplicate ticketing overhead, and a mean end-to-end processing latency of 18.4\,ms per complaint. The complete system guarantees strict identity confidentiality while dramatically accelerating civic response velocity.
\end{abstract}

\begin{IEEEkeywords}
Anonymous Systems, Multi-Agent Systems, Natural Language Processing, Civic Computing, Personally Identifiable Information (PII) De-identification, Semantic Triage, Duplicate Detection.
\end{IEEEkeywords}

\section{Introduction}
\label{sec:intro}
Public grievance redressal platforms serve as the indispensable interface between citizens and municipal administrations \cite{worldbank2020governance}. Conventional electronic governance systems---such as India's Centralised Public Grievance Redress and Monitoring System (CPGRAMS) \cite{cpgrams2022} or municipal web portals---mandate exhaustive identity disclosures including citizens' legal names, verified telephone numbers, residential postal addresses, and national identity credentials (e.g., Aadhaar) prior to complaint registration. While conceived to suppress frivolous submissions, this mandatory identity disclosure induces a profound \textit{civic trust deficit} and creates an acute \textit{whistleblower paradox} \cite{whistleblower2021}. Citizens encountering institutional corruption, hazardous municipal negligence, or organized extortion routinely withhold critical civic intelligence due to legitimate apprehensions of official harassment, physical intimidation, or socio-economic retribution.

Concurrently, public administration departments are overwhelmed by the velocity and volume of citizen submissions \cite{almarabeh2016electronic}. Existing triage pipelines rely heavily on human administrative staff who manually read incoming tickets, guess departmental allocations, assign subjective urgency labels, and attempt to cross-reference existing complaints. This manual workflow exhibits three critical failure modes:
\begin{enumerate}
    \item \textbf{Identity Leakage and Traceability:} Complainant identifiers are exposed across intermediate clerical hierarchies, creating pervasive privacy vulnerabilities and deterring whistleblowers.
    \item \textbf{Cognitive Bottlenecks in Urgency Triage:} Life-threatening civic emergencies (e.g., active high-voltage electrical sparking, open highway sewer trenches, toxic drinking water contamination) languish in chronological queues behind low-severity cosmetic defects (e.g., burned-out streetlights).
    \item \textbf{Duplicate Congestion and Administrative Paralysis:} Localized infrastructure failures routinely trigger hundreds of identical complaints from proximate residents, inundating departmental queues with redundant records and obscuring novel emergencies.
\end{enumerate}

To decisively resolve these systemic vulnerabilities, this paper introduces \textbf{JanAwaaz AI} (\textit{`Voice of the People''}), an autonomous multi-agent computing framework that reconciles absolute citizen anonymity with high-throughput, deterministic civic triage. By decoupling identity verification from grievance tracking via cryptographic one-way hashes and deploying an asynchronous pipeline of specialized natural language processing (NLP) agents, JanAwaaz AI ensures that no personally identifiable information (PII) is ever committed to persistent storage or exposed to municipal personnel.

\subsection{Key Contributions}
The core scientific and architectural contributions of this work are summarized as follows:
\begin{itemize}
    \item \textbf{Privacy-Preserving Ingestion Architecture:} We design a linear-time, zero-knowledge PII de-identification agent combining regular expression automata and contextual named-entity filters to systematically sanitize contact numbers, emails, national identity tokens, and self-identifying markers before persistent database serialization.
    \item \textbf{Cryptographic Tracking Without Authentication:} We introduce an anonymous transaction tracking protocol that generates collision-resistant SHA-256 tokens ({\text{track}}$) parameterized on complaint entropy, enabling citizens to monitor resolution lifecycle states without establishing user accounts, cookies, or digital footprints.
    \item \textbf{Multi-Factor Heuristic Urgency Scoring Engine:} We formulate an empirical scoring metric ($\Phi_{\text{urgency}}$) integrating safety-critical domain lexicons, syntactic panic indicators (exclamation density), lexical capitalization ratios, and textual length heuristics to categorize complaints into four deterministic triage tiers: \textit{Low}, \textit{Medium}, \textit{High}, and \textit{Critical}.
    \item \textbf{Vector-Space Duplicate Clustering Agent:} We deploy an online term-frequency inverse-document-frequency (TF-IDF) vectorizer with cosine similarity thresholding ($\theta_{\text{dup}} = 0.55$) that dynamically identifies and links duplicate complaints to root records in real time, preventing queue flooding while preserving demographic complaint frequency.
    \item \textbf{Production Cloud Architecture \& Empirical Validation:} We implement the full stack using FastAPI asynchronous event loops, SQLAlchemy ORM, React 18 Progressive Web Application (PWA), and React Native Expo mobile architecture. We deploy the system on Render cloud microservices and benchmark it across 500+ diverse civic scenarios, demonstrating sub-20\,ms triage latency, 99.4\% PII stripping precision, and 91.8\% routing accuracy.
\end{itemize}
